\section{Multiple Roots}

\begin{definition}
	Let \( K \) be a field, \( 0\ne f \in K[x] \), recall that \( f(a) = 0\iff (x-a) \mid f \). \( a
	\in K \) is a root of \( f \) with multiplicity \( m \) if \( (x-a)^m \mid f \) and \( (x-a)^{m+1} \nmid f \).
	\begin{itemize}
		\item Simple root: root of multiplicity \( 1 \).
		\item Multiple root: root of multiplicity \( \ge 2 \).
	\end{itemize}
\end{definition}
To check whether it is a simple root, one can use \textbf{derivative}.

\begin{definition}
	If
	\[
		f = a_n x^n + a_{n-1} x^{n-1} + \cdots + a_2 x^2 + a_1 x + a_0 \in K[x], \quad a_n \ne 0
	\]
	then define
	\[
		f' = n a_n x^{n-1} + (n-1) a_{n-1} x^{n-2} + \cdots + 2 a_2 x + a_1 \in K[x].
	\]
\end{definition}

\begin{remark}[\textbf{Properties}]
	\leavevmode
	\begin{itemize}
		\item \( f \to f' \) is \( K \)-linear.
		\item (Leibniz Rule) \( (fg)' = fg' + f'g \).
	\end{itemize}
\end{remark}

\begin{proof}
	\leavevmode
	\begin{itemize}
		\item 1 is easy.
		\item For 2, consider the following map
		      \[
			      (f,g) \mapsto (fg)' - fg' - f'g
		      \]
		      is \( K \)-linear in both \( f,g \). Hence it is enough to check it when \( f = x^i,g=x^j \).
		      We have
		      \[
			      (fg)' = (x^{i+j})' = (i+j)x^{i+j-1}
		      \]
		      and we have
		      \[
			      fg' + f'g = x^i (jx^{j-1}) + (ix^{i-1})x^j = (fg)'
		      \]
	\end{itemize}
\end{proof}

\begin{proposition}
	If \( 0 \ne f \in K[x] \) and \( a \in K \). \( a \) is a multiple root of \( f \) iff \( f(a) =
	0, f'(a) = 0 \).
\end{proposition}

\begin{proof}
	\( f(a) = 0 \) iff \( (x-a) \mid f \). In this case, we can decompose \( f \)
	\[
		f = (x-a) g \text{ for some } g \in K[x]
	\]
	then
	\[
		\begin{aligned}
			f'            & = g + (x-a)g'                                      \\
			f'(a)         & = g(a)                                             \\
			\implies f(a) & = 0 = f'(a)                                        \\
			\iff f        & = (x-a)g \text{ for some } g \text{ and } g(a) = 0 \\
			\iff (x-a)^2  & \mid f
		\end{aligned}
	\]
\end{proof}

\section{Classification of Finite Fields}
Suppose that \( K \) is a \textcolor{blue}{finite field}, then \( \text{char}(K) \ne 0 \). Let \( p \) be some
prime such that \( \text{char}(K) = p \). Then there exists some \( \FF_p = \qo{\ZZ}{p\ZZ} \), such
that
\[
	\FF_p \hookrightarrow K
\]
and let \( e = [K: \FF_p] < \infty \), we have \( \abs{K} = p^e \). Thus \( \abs{K^\times} = p^e - 1 \). By Lagrange theorem,
\( a^{p^e - 1} = 1 \; \forall \; a\in K, a\ne 0 \). Thus \( a^{p^e} = a, \forall \; a\in K
\smallsetminus \{0\} \).

\begin{theorem}
	\leavevmode
	\begin{enumerate}
		\item For every \( e \geq 1 \), \( \exists \) field with \( p^e \) elements.
		\item Every two such fields are isomorphic.
	\end{enumerate}
\end{theorem}

\begin{proof}
	Let \( \overline{\FF_p} \) be an algebraic closure of \( \FF_p \). Let
	\[
		K = \{a\in \overline{\FF_p} \; | \; a^{p^e} - a= 0\}
	\]
	\begin{claim}
		\( K \subseteq \overline{\FF_p} \) is a subfield.
	\end{claim}

	\begin{proof}
		It is clear that \( 0,1\in K \) and \( a,b\in K \implies ab \in K \). We saw the frobenius
		homomorphism \ref{prop:frobenius-hom} that:
		\[
			F: \overline{\FF_p} \to \overline{\FF_p}, \quad F(a) = a^p
		\]
		is a ring homomorphism. Thus \( F^e \) is also a ring homomorphism. So
		\[
			\{a \in \overline{\FF_p} \; | \; F^e(a) = a\} \subseteq \overline{\FF_p} \text{ is subgroup.}
		\]
		\( \implies \)
		\[
			K \subseteq \overline{\FF_p} \text{ is subring. } (\implies \ne \{0\}, \text{ commutative.})
		\]
		See that \( a \in K^\times \implies a^{p^e-1} = 1 \implies a^{-1} = a^{p^e - 2} \in K \implies K \) is a field.
		\( K \) has \( p^e \) elements: Let \( f = x^{p^e} - x \) decompose in \( \overline{\FF_p}[x] \)
		as
		\[
			\begin{aligned}
				f  & = x^{p^e} -x = \prod_{i=1}^{p^e} (x-\alpha_i)                     \\
				f' & = p^e x^{p^e-1} - 1 = -1 \text{ with } p^e \equiv 0 \; (\bmod{p})
			\end{aligned}
		\]
		with \( f' = -1 \implies \) all roots of \( f \) are simple \( \implies \abs{K} = p^e \). This
		follows from the definition of \( K \), which are the elements in \( \overline{\FF_p} \) that
		makes the roots of polynomial \( f \).
	\end{proof}
	\textbf{Uniqueness:} Suppose \( L \) be a field with \( p^e \) elements. Have
	\[
		\begin{tikzcd}
			& L \arrow[dashed, rdd, "\vp" ] & \\
			\FF_p \arrow[ru, hook] \arrow[rd, hook] & \\
			& K \arrow[r, hook] & \overline{\FF_p}
		\end{tikzcd}
	\]
	with \( \qo{L}{\FF_p} \) being algebraic since it is finite, thus by \textbf{Proposition}
	\ref{prop:extension-alg}, there exists a morphism of \( \FF_p \)-algebra \( \vp: L \to
	\overline{\FF_p} \). Since \( a^{p^e} = a \; \forall \; a\in L \) (since \( \abs{L^\times} = p^e -
	1\)), then \( \vp(a)^{p^e} = \vp(a) \; \forall \; a\in L \implies \vp: L \hookrightarrow K
	\implies \abs{L} = \abs{K} \) and the fact that \( \vp \) is injective, this is isomorphism.
\end{proof}

\section{Splitting Fields}

\begin{definition}[\textbf{Splitting Fields}]
	Let \( K \) be a field, \( f\in K[x] \) of degree \( \ge 1 \). A \underline{splitting field} of \( f \) over
	\( K \) is a field extension \( K \hookrightarrow L \), such that
	\begin{enumerate}
		\item \( f \) splits as a product of degree \( 1 \) polynomial in \( L[x] \).
		\item This does not hold for every \( K \subseteq L' \subsetneq L \) with \( L' \) be a
		      subextension of \( L \).
	\end{enumerate}
	\textbf{Equivalently}, we have
	\[
		f = c(x-\alpha_1)\cdots (x-\alpha_n)
	\]
	and
	\[
		\textcolor{blue}{L = K(\alpha_1,\ldots,\alpha_n)}
	\]
\end{definition}

\begin{eg}
	Let \( f = x^3 - 2 \). Its roots in \( \CC \) are \( \sqrt[3]{2}, \sqrt[3]{2}\omega, \sqrt[3]{2}\omega^2 \) where
	\( \omega = \cos \frac{2\pi}{3} + i\sin\frac{2\pi}{3} \). Here a splitting field for \( f \) over
	\( \QQ \) is \( \QQ(\sqrt[3]{2}, \omega) \).
\end{eg}

\begin{proposition}
	If \( f \in K[x] \) has degree \( \ge 1 \), then
	\begin{enumerate}
		\item \( f \) has a splitting field over \( K \).
		\item Any two such splitting field are isomorphic as extensions of \( K \).
	\end{enumerate}
\end{proposition}

\begin{proof}
	\leavevmode
	\begin{enumerate}
		\item Consider an algerbaic closure \( \overline{K} \) of \( K\).
		      \[
			      f = c (x-\alpha_1)\cdots (x-\alpha _n) \quad \alpha _1,\ldots,\alpha_n \in \overline{K},
			      c\in K
		      \]
		      See that \( K(\alpha_1, \ldots,\alpha_n) \) is a splitting field of \( f \) over \( K \) and
		      yields the existence.
		\item Suppose \( L \) is another splitting field of \( f \) over \( K \), we have
		      \[
			      f = c (x-\beta_1)\cdots (x-\beta_n) \quad \beta_1,\ldots,\beta_n \in L
		      \]
		      thus we have
		      \[
			      \begin{tikzcd}
				      & L = K(\beta_1,\ldots, \beta_n) \arrow[dashed, rdd, "\vp" ] & \\
				      K \arrow[ru,hook] \arrow[rd,hook ] & & \\
				      & K(\alpha_1,\ldots, \alpha_n) \arrow[r,hook] &
				      \overline{K}
			      \end{tikzcd}
		      \]
		      with \( \qo{L}{K} \) algebraic since it is genrated by a bunch of algerbaic elements. By
		      \textbf{Proposition} \ref{prop:extension-alg}, \( \vp \) exists. See that
		      \( f(\beta_i) = 0 \; \forall \; i \implies f(\vp(\beta_i)) = 0 \implies \vp(\beta_i) \in \{\alpha_1,\ldots,\alpha_n\} \).
		      \[
			      \implies \vp : L \hookrightarrow K(\alpha_1,\ldots,\alpha_n)
		      \]
		      This is isomorphism since \( f \) can't be split over any proper subextension of \(
		      K(\alpha_1,\ldots,\alpha_n) \) by the definition of splitting field.
	\end{enumerate}
\end{proof}

\section{Separable Extensions}
Let \( K \hookrightarrow L \) be field extension, \( a\in L \) is algerbaic over \( K \) with \(
f\in K[x] \) be its minimal polynomial. Consider \( q = \gcd(f,f') \) with \( f' \) be the
derivative of \( f \). \( f \) is the minimal polynomial meaning that \( f \) is irreducible in \(
K[x] \). Then there are two cases for \( q \):
\begin{itemize}
	\item \( q \sim f \implies f\mid f' \), but see that \( \deg(f') < \deg(f) \implies f' = 0 \).
	      This holds iff \( \text{char}(K) = p >0 \) and \( f \in K[x^p] \). In this case, every root of \( f \) is a multiple root (\( f' = 0 \)).
	\item \( q \sim 1 \implies \exists \; u,v \in K[x] \), such that
	      \[
		      uf + vf' = 1
	      \]
	      \( \implies f, f' \) have no common root in any extension of \( K \), otherwise contradict to
	      the above equation (\( 0=1 \)), thus all roots of \( f \) are simple.
\end{itemize}

\begin{definition}[\textbf{Separable}]
	\leavevmode
	\begin{enumerate}
		\item \( a\in L \) is \underline{separable} over \( K \) if \( a \) is a simple root of
		      its minimal polynomial over \( K \).
		\item \( \qo{L}{K} \) algerbaic is separable if \( \forall \; a \in L \), \( a \) is separable
		      over \( K \).
	\end{enumerate}
\end{definition}

